\section{Restricted survivors and the spectral mechanism}\label{sec:survivors}
\subsection{A uniform qubit bound for relative variation}
The failure of uniform concentration in dimension four coexists with a
restricted theorem for qubits.

\begin{proposition}[Qubit coefficient-of-variation bound]\label{prop:qubit-cv}
For a nonstationary qubit under a time-independent Hermitian Hamiltonian,
extend $r,q,R$ through their removable recurrence values. For any
probability-weighted time average,
\begin{equation}\label{eq:qubit-cv}
 \CV(r),\ \CV(q),\ \CV(R)
 \ \le\ \frac{\kappa_*-1}{2\sqrt{\kappa_*}},
 \qquad \kappa_*=\frac{17+7\sqrt7}{27}.
\end{equation}
The upper bound is approximately $0.1375633885$. No claim is made that
this value is attained by a qubit trajectory and averaging measure.
\end{proposition}
\begin{proof}
Put $y=2\sqrt{p_1p_2}\in[0,1)$ and $h=\sin^2\theta>0$; nonstationarity
also requires a nonzero energy gap. The coefficients in
Eq.~\eqref{eq:qubit-mus} become
\begin{equation}\label{eq:qubit-range}
 \mu_S=\frac{(1-y)h}{2},\qquad
 \mu_K=\frac{(1-y^2)h}{2-y^2},\qquad
 \kappa=\frac{\max_t r(t)}{\min_t r(t)}
 =\frac{1-\mu_S}{1-\mu_K}.
\end{equation}
The same range factor applies to $q$ and $R$. It increases with $h$, so
\begin{equation}\label{eq:kappa-max}
 \kappa\le\frac{(1+y)(2-y^2)}2
 \le\frac{17+7\sqrt7}{27},\qquad
 y_{\max}=\frac{\sqrt7-1}{3}.
\end{equation}
The second derivative of the middle expression is $-1-3y<0$, and its
endpoint values are one. For a positive random variable $X\in[m,M]$,
the inequality $\mathbb E[(X-m)(M-X)]\ge0$ implies
\begin{equation}\label{eq:bounded-variance}
 \Var(X)\le(M-\mathbb EX)(\mathbb EX-m),\qquad
 \CV(X)\le\frac{M-m}{2\sqrt{mM}}
 =\frac{\kappa-1}{2\sqrt\kappa}.
\end{equation}
The last step maximizes the relative variance over the allowed mean.
Apply this bound to each of the three ratios.
\end{proof}

Completely stationary $0/0$ ratios are excluded from this proposition.
Dimension four suffices for the temporal no-go constructions, while the
qubit bound excludes dimension two. Whether dimension three suffices for
unbounded temporal variation is not settled here. No minimum-dimension
claim is made for the fixed-purity no-go theorem.

\subsection{State-dependent comparisons}
The return bound yields useful comparisons retaining the two spectral
measures. Let $m_S,m_K$ be the finite lengths of the $\SU$ and $\CM$
chains and define
\begin{equation}\label{eq:returns}
 A_S=\Tr\bigl(\sqrt{\rho(t)}\sqrt\rho\bigr),\qquad
 A_K=\frac{\Tr(\rho(t)\rho)}P.
\end{equation}
\begin{corollary}[Return comparison]\label{cor:return-comparison}
When both return losses are nonzero,
\begin{equation}\label{eq:state-dependent}
 \frac{1-A_S^2}{(m_K-1)(1-A_K^2)}
 \le r(t)\le
 \frac{(m_S-1)(1-A_S^2)}{1-A_K^2}.
\end{equation}
\end{corollary}
\begin{proof}
Divide the lower return bound for the numerator by the upper one for
the denominator, and conversely.
\end{proof}

These bounds involve return amplitudes and cyclic dimensions, rather
than purity alone. Similarly, a necessary short-time condition for the
source hierarchy would be
\begin{equation}\label{eq:curvature-necessary}
 \norm{[H,\sqrt\rho]}_{\HS}^2
 \le\frac{\norm{[H,\rho]}_{\HS}^2}{P}
 \le2\Var_\rho(H),\qquad
 \Var_\rho(H)=\Tr(\rho H^2)-(\Tr\rho H)^2.
\end{equation}
The displayed counterexamples violate the respective comparisons.
Satisfying these necessary curvature inequalities is not claimed to be
sufficient for an all-time hierarchy.

\subsection{Why the comparisons separate}
The mixed seed and the square-root seed give different powers of a
population weight to a coherent sector. Hilbert--Schmidt normalization
then introduces the global purity into only the former. A highly populated
stationary sector can suppress mixed-state operator complexity more strongly
than purification spread complexity. Independent block frequencies permit
partial recurrences that amplify this mismatch in time. Holding the total
purity fixed does not fix the distribution of population among these sectors.

By comparison, a pure density matrix is a tensor product after
vectorization. Its difference-convolution measure comes with the nested
polynomial subspaces of Section~\ref{sec:factor}. No corresponding
containment is generally available for the $\rho$ and $\sqrt\rho$
measures. Qubits are protected by their common three-atom form and the
monotone ordering of the three coefficients. The loss of these structures,
rather than an inconsistency between state and operator definitions,
explains the different outcomes.
